\section*{第 \ref{ch:g11:gene-expression} 章——从基因到蛋白质}

\begin{solution}{exo:g11:gene-expression:1}
\omterm{def:g11:gene-expression:transcription}{RNA} 为单链，其糖为核糖，并以尿嘧啶代替胸腺嘧啶
（它也寿命较短，且会离开\omterm{def:g10:cells-common-unit:organelle}{细胞核}）。
\end{solution}

\begin{solution}{exo:g11:gene-expression:2}
\texttt{5'-AUG GCU UUU UAA-3'}.
\end{solution}

\begin{solution}{exo:g11:gene-expression:3}
Met--Gly--Lys--Cys，然后终止。
\end{solution}

\begin{solution}{exo:g11:gene-expression:4}
\omterm{def:g10:cells-common-unit:organelle}{核糖体}逐个\omterm{def:g11:gene-expression:code}{密码子}地读取 mRNA，把\omterm{def:g10:chemistry-of-life:families}{氨基酸}连接成肽链；
每种转运 \omterm{def:g11:gene-expression:transcription}{RNA} 借其反\omterm{def:g11:gene-expression:code}{密码子}与一个\omterm{def:g11:gene-expression:code}{密码子}配对，
并携带与之对应的\omterm{def:g10:chemistry-of-life:families}{氨基酸}。
\end{solution}

\begin{solution}{exo:g11:gene-expression:5}
若用 4 个字母，两字母一组可得 $4^2 = 16$ 种组合，
少于 20 种\omterm{def:g10:chemistry-of-life:families}{氨基酸}；三字母一组可得 64 种，足够。
\end{solution}

\begin{solution}{exo:g11:gene-expression:6}
$412 \times 3 + 3 = 1239$ 个\omterm{def:g10:universal-dna:nucleotide}{核苷酸}；\omterm{def:g10:universal-dna:gene}{基因}至少需含 1239 个
碱基对（真核生物因含\omterm{prop:g11:gene-expression:splicing}{内含子}，数目更多）。
\end{solution}

\begin{solution}{exo:g11:gene-expression:7}
原序列：Met--Lys--Gly--Trp。\omterm{def:g11:mutations:mutation}{突变}体：UGA 为终止\omterm{def:g11:gene-expression:code}{密码子}，故只得 Met--Lys--Gly：
无义\omterm{def:g11:mutations:mutation}{突变}，使\omterm{def:g11:gene-expression:protein}{蛋白质}提前截短。
\end{solution}

\begin{solution}{exo:g11:gene-expression:8}
\texttt{AAG} 仍编码 Lys：属同义\omterm{def:g11:mutations:mutation}{突变}，Met--Lys--Gly--Trp 不变。
第二种是插入一个 G，使阅读框移动：\texttt{AUG AAA GGG
CUG GUA A}\dots{} Met--Lys--Gly--Leu--Val\dots，为移码\omterm{def:g11:mutations:mutation}{突变}，
该片段内无终止\omterm{def:g11:gene-expression:code}{密码子}。
\end{solution}

\begin{solution}{exo:g11:gene-expression:9}
密码表的多数区块中，前两个字母相同的四个\omterm{def:g11:gene-expression:code}{密码子}
编码同一种\omterm{def:g10:chemistry-of-life:families}{氨基酸}（Pro、Thr、Ala、Gly、Val、Ser、
Leu\dots）：改变第三个字母则不改变所编码的\omterm{def:g10:chemistry-of-life:families}{氨基酸}。
简并度主要集中在第三位。
\end{solution}

\begin{solution}{exo:g11:gene-expression:10}
按三个一组读取，\texttt{UCU CUC UCU CUC}\dots{} 交替出现两个\omterm{def:g11:gene-expression:code}{密码子}，
即 UCU 与 CUC，因此\omterm{def:g11:gene-expression:protein}{蛋白质}交替含两种\omterm{def:g10:chemistry-of-life:families}{氨基酸}，与观察一致。若为两字母码，
同样的重复只会对应一个\omterm{def:g11:gene-expression:code}{密码子}和一种\omterm{def:g10:chemistry-of-life:families}{氨基酸}。
由于多聚 U 编码 Phe，而表中 UCU 为 Ser、CUC 为 Leu，
此实验应将 UCU 指定为 Ser、CUC 指定为 Leu（或相反，
须由其他实验确定）。
\end{solution}

\begin{solution}{exo:g11:gene-expression:11}
编码序列仅需 1503 个碱基对；
其余部分为\omterm{prop:g11:gene-expression:splicing}{内含子}，\omterm{def:g11:gene-expression:transcription}{转录}后从前体 mRNA 中剪除，
另外\omterm{def:g10:universal-dna:gene}{基因}两端还有调控序列。
\end{solution}

\begin{solution}{exo:g11:gene-expression:12}
阻断\omterm{prop:g11:gene-expression:translation}{翻译}可杀死细菌，同时保持患者\omterm{def:g10:cells-common-unit:cell}{细胞}正常工作。
该药物的选择性表明，细菌与人类的
\omterm{def:g10:cells-common-unit:organelle}{核糖体}虽然以相同密码完成相同工作，但结构差异
足以使某一分子只与其中一种结合——
这是两者自分化以来积累的差异。
\end{solution}

\begin{solution}{exo:g11:gene-expression:13}
小鼠的 \omterm{def:g11:gene-expression:transcription}{RNA} 聚合酶\omterm{def:g11:gene-expression:transcription}{转录}了水母\omterm{def:g10:universal-dna:gene}{基因}，其\omterm{def:g10:cells-common-unit:organelle}{核糖体}与 tRNA
使用相同的密码表\omterm{prop:g11:gene-expression:translation}{翻译} mRNA，因此
生成了相同的\omterm{def:g10:chemistry-of-life:families}{氨基酸}序列，即相同的荧光蛋白。
若密码不同，\omterm{def:g11:gene-expression:code}{密码子}将被读为其他\omterm{def:g10:chemistry-of-life:families}{氨基酸}，
\omterm{def:g11:gene-expression:protein}{蛋白质}便成为无意义、不发荧光的肽链。
\end{solution}

\begin{solution}{exo:g11:gene-expression:14}
UAG 将不再终止\omterm{prop:g11:gene-expression:translation}{翻译}：凡\omterm{def:g10:universal-dna:gene}{基因}以 UAG 结尾的\omterm{def:g11:gene-expression:protein}{蛋白质}，
都会延长至正常末端之后，常常丧失功能——
这是一种普遍的危害。但若\omterm{def:g10:universal-dna:gene}{基因}因提前出现的 UAG 而截短，
现在则会被读通，恢复接近全长的\omterm{def:g11:gene-expression:protein}{蛋白质}：
第二个\omterm{def:g11:mutations:mutation}{突变}抑制了第一个\omterm{def:g11:mutations:mutation}{突变}。
\end{solution}

\begin{solution}{exo:g11:gene-expression:15}
移码\omterm{def:g11:mutations:mutation}{突变}会打乱其下游的全部编码，通常在数个\omterm{def:g11:gene-expression:code}{密码子}之内
即遇到终止\omterm{def:g11:gene-expression:code}{密码子}：若发生在起始端附近，几乎得不到正确的\omterm{def:g11:gene-expression:protein}{蛋白质}；
若发生在末端附近，\omterm{def:g11:gene-expression:protein}{蛋白质}大部分完整，仍可能折叠并发挥作用。
错义\omterm{def:g11:mutations:mutation}{突变}只改变一个\omterm{def:g10:chemistry-of-life:families}{氨基酸}：若该位置能耐受改变则无害，
若位于活性位点或破坏了折叠结构则是致命的——
效应取决于改变的是哪一个\omterm{def:g10:chemistry-of-life:families}{氨基酸}、位于何处。
\end{solution}

\begin{solution}{pb:g11:gene-expression:1}
\textbf{1.} \texttt{3'-TAC CGA ACC TTT GGG CTT ATG CCA GTA AAG ACT-5'}.

\textbf{2.} \texttt{AUG GCU UGG AAA CCC GAA UAC GGU CAU UUC UGA}.

\textbf{3.} Met--Ala--Trp--Lys--Pro--Glu--Tyr--Gly--His--Phe，终止。

\textbf{4.} 10 个\omterm{def:g10:chemistry-of-life:families}{氨基酸}；33 个\omterm{def:g10:universal-dna:nucleotide}{核苷酸}（共 11 个\omterm{def:g11:gene-expression:code}{密码子}，
含终止\omterm{def:g11:gene-expression:code}{密码子}）。

\textbf{5.} AUG；第一个 AUG 的位置决定阅读框，
此后的每个\omterm{def:g11:gene-expression:code}{密码子}都自它开始依次读取。

\textbf{6.} CCG 仍编码 Pro：同义\omterm{def:g11:mutations:mutation}{突变}，\omterm{def:g11:gene-expression:protein}{蛋白质}不变。

\textbf{7.} CGU 编码 Arg 替代 His：错义\omterm{def:g11:mutations:mutation}{突变}，
Met--Ala--Trp--Lys--Pro--Glu--Tyr--Gly--Arg--Phe.

\textbf{8.} UGA 为终止\omterm{def:g11:gene-expression:code}{密码子}：Met--Ala，无义\omterm{def:g11:mutations:mutation}{突变}。仅
第三个\omterm{def:g11:gene-expression:code}{密码子}中一个 G 替换为 A（TGG 变为 TGA）。

\textbf{9.} mRNA \texttt{AUG GCU UUG GAA ACC CGA AUA CGG UCA UUU CUG
A}：Met--Ala--Leu--Glu--Thr--Arg--Ile--Arg--Ser--Phe--Leu\dots，片段中无终止\omterm{def:g11:gene-expression:code}{密码子}，
移码\omterm{def:g11:mutations:mutation}{突变}：第二个\omterm{def:g10:chemistry-of-life:families}{氨基酸}之后的每个\omterm{def:g10:chemistry-of-life:families}{氨基酸}都
发生改变。

\textbf{10.} \omterm{def:g11:mutations:mutation}{突变}体 1（同义\omterm{def:g11:mutations:mutation}{突变}，无变化）$<$ \omterm{def:g11:mutations:mutation}{突变}体 2（一个\omterm{def:g10:chemistry-of-life:families}{氨基酸}
改变，功能可能保留）$<$ \omterm{def:g11:mutations:mutation}{突变}体 4（移码\omterm{def:g11:mutations:mutation}{突变}，除前两个外
\omterm{def:g10:chemistry-of-life:families}{氨基酸}全部错误）$\approx$ \omterm{def:g11:mutations:mutation}{突变}体 3（无义\omterm{def:g11:mutations:mutation}{突变}，两个\omterm{def:g10:chemistry-of-life:families}{氨基酸}的
片段，\omterm{def:g11:gene-expression:protein}{蛋白质}已丧失）。

\textbf{11.} Leu 6 个，Ser 6 个，Trp 1 个。Trp 仅有一个\omterm{def:g11:gene-expression:code}{密码子}，
任何替换都会使其改变；Leu 或 Ser 有六个\omterm{def:g11:gene-expression:code}{密码子}，多数情况下不变。

\textbf{12.} CUU、CUC、CUA、CUG 均编码 Leu：第三位
的 3 种可能替换全部为同义\omterm{def:g11:mutations:mutation}{突变}——100\%。

\textbf{13.} 两个字母只能给出 16 个\omterm{def:g11:gene-expression:code}{密码子}，不足以编码 20 种\omterm{def:g10:chemistry-of-life:families}{氨基酸}
加上终止\omterm{def:g11:gene-expression:code}{密码子}；三个字母给出 64 个，已足够，因此四个字母（256 个）
便是浪费。

\textbf{14.} UUU $=$ Phe，AAA $=$ Lys。

\textbf{15.} \omterm{def:g10:universal-dna:gene}{基因}在\omterm{def:g10:biodiversity-scales:species}{物种}间转移后，读取方式完全相同，
因此\omterm{prop:g10:universal-dna:universal}{转基因技术}可行；通用的\omterm{def:g11:gene-expression:code}{遗传密码}一代代保持不变，
是所有生物体源自共同祖先的证据，该祖先
早已使用这套密码。

\textbf{16.} 10 个\omterm{def:g10:chemistry-of-life:families}{氨基酸}需 \qty{1}{s}，600 个\omterm{def:g10:chemistry-of-life:families}{氨基酸}需 \qty{60}{s}。

\textbf{17.} $1800/80 \approx 22$ 个\omterm{def:g10:cells-common-unit:organelle}{核糖体}同时工作。

\textbf{18.} $7200/8 = 900$ 个拷贝。

\textbf{19.} 一个 mRNA 在 10 分钟内产生 $600/8 = 75$ 个拷贝；
约需 27 个 mRNA。

\textbf{20.} 该\omterm{def:g11:gene-expression:protein}{蛋白质}为
Met--Ala--Trp--Lys--Pro--Glu--Tyr--Gly--His--Phe；\omterm{def:g11:mutations:mutation}{突变}体 3，
单个 G 变为 A 的替换在第三个\omterm{def:g11:gene-expression:code}{密码子}处产生终止，使其丧失；
一个\omterm{def:g10:cells-common-unit:organelle}{核糖体}约在一秒内合成这 10 个\omterm{def:g10:chemistry-of-life:families}{氨基酸}。
\end{solution}
